%% Example CV - Alessandro Chiodo
%% Comprehensive overview of all competencies, experience, and achievements.
%% Use as a master reference when tailoring role-specific CVs.
%%
%% Compile with: lualatex main_example.tex
%% (pdflatex often fails on modern MiKTeX with fontawesome5 font-expansion errors.)

\documentclass[11pt,a4paper,sans]{moderncv}
\moderncvstyle{banking}
\moderncvcolor{blue}

% Force the name and section headings to render in moderncv blue (color1).
% Default banking leaves them black: moderncvstylebanking.sty's \colorlet
% copies (not aliases) the pre-scheme accent colour, so the name colours are
% frozen before \moderncvcolor runs. Re-let them after. \namefont is the hook
% every name-style macro routes through, so this also works on moderncv 2.3.1
% (Debian/Ubuntu apt), which has no \firstnamestyle/\lastnamestyle at all.
\renewcommand*{\namefont}{\fontsize{34}{36}\bfseries\upshape}
\colorlet{firstnamecolor}{color1}
\colorlet{lastnamecolor}{color1}
\colorlet{namecolor}{color1}
\renewcommand*{\sectionstyle}[1]{{\sectionfont\color{color1}#1}}

\usepackage[utf8]{inputenc}
% moderncv loads hyperref itself in an \AtEndPreamble hook, so \hypersetup
% must go in an \AtEndPreamble of our own: on moderncv < 2.4 a top-level
% \usepackage{hyperref} clashes with the class's own
% \RequirePackage[unicode]{hyperref}. From 2.4.0 the class passes its options
% through \PassOptionsToPackage instead, which is what removes that clash.
\AtEndPreamble{\hypersetup{
    colorlinks=true,
    linkcolor=blue,
    filecolor=magenta,
    urlcolor=blue,
    pdftitle={Alessandro Chiodo - CV},
    % Keep pdfpagemode=UseNone: this block runs after moderncv's own
    % \AtEndPreamble (moderncv.cls sets pdfpagemode there), so a FullScreen
    % value here would win and open every CV in fullscreen presentation mode.
    pdfpagemode=UseNone,
}}
\usepackage[scale=0.80]{geometry}
\usepackage{import}

% personal data
\name{Alessandro}{Chiodo}
\address{Palermo, Italia (Disponibile Full Remote)}{}{}
\phone[mobile]{+39 351 7727138}
\email{alexbypa@gmail.com}

\begin{document}

\makecvtitle

% ============================================================
%     PROFILE STATEMENT
% ============================================================

\vspace{6pt}
\small{Senior Software Developer con oltre 20 anni di esperienza nello sviluppo di servizi backend, applicazioni web e integrazioni enterprise. Competenze consolidate in C\#, .NET, SQL Server, REST API e messaggistica asincrona, maturate in sistemi ad alta affidabilità e contesti regolamentati. Orientato a performance, manutenibilità, architetture a basso accoppiamento e qualità del software; esperienza recente nell’applicazione di AI, LLM e Model Context Protocol a sistemi .NET. Appartenente alle Categorie Protette ai sensi della Legge 68/99, Iscritto al Collocamento Mirato dal 28/07/2026.}

% ============================================================
%     CORE COMPETENCIES
% ============================================================

\section{Competenze Tecniche}
\vspace{1pt}
\begin{itemize}

\item \textbf{Backend e API}: C\#, .NET, ASP.NET, REST API, LINQ, Dependency Injection, Design Patterns, Clean Architecture.

\item \textbf{Database e Performance}: SQL Server, PostgreSQL, Redis; Stored Procedure, trigger, CTE, tuning query e gestione memoria.

\item \textbf{Integrazioni e Asincronia}: RabbitMQ, Kafka, HTTP, SOAP/XML, System.Threading.Channels, Worker Services.

\item \textbf{Testing e Delivery}: xUnit, Moq, Postman, Newman, Git; utilizzo operativo di Azure DevOps e CI/CD.

\item \textbf{Osservabilità}: Logging strutturato, Elasticsearch, OpenTelemetry, Serilog, Microsoft.Extensions.Logging.

\item \textbf{AI Engineering}: MCP (Model Context Protocol), Microsoft.Extensions.AI, Semantic Kernel, Ollama, RAG, SQL guardrails per agenti AI.

\end{itemize}

% ============================================================
%     PROFESSIONAL EXPERIENCE
% ============================================================

\section{Esperienza Professionale}
\vspace{3pt}
\begin{itemize}

% --- Pixelo ---
\item{\cventry{10/2017--11/2025}{Sviluppatore capo progetto}{Pixelo S.r.l.}{Palermo, Italia}{}{\vspace{1pt}
\begin{itemize}
    \item Progettazione di REST Web Services in C\# con Dependency Injection e Design Patterns in contesto regolamentato ADM/AAMS; centralizzazione della logica in un hub core condiviso, riducendo di circa 25 volte i tempi di sviluppo dei nuovi servizi.
    \item Gestione e ottimizzazione di SQL Server (Stored Procedure, trigger e CTE) per circa 300 transazioni concorrenti, con tempo medio di elaborazione di 7-8 ms.
    \item Implementazione di integrazioni asincrone RabbitMQ e Kafka con pattern proxy/factory configurabile, per elaborare messaggi di fornitori con modelli dati eterogenei.
    \item Sviluppo dell'estensione Azure DevOps NugetHelper in Node.js e integrazione nelle pipeline CI/CD per automatizzare build, distribuzione di artefatti e deployment simultaneo di circa 25 servizi.
    \item Realizzazione di un pacchetto interno di logging con Elasticsearch e OpenTelemetry, distribuito tramite Azure DevOps Artifacts.
\end{itemize}}}

\vspace{3pt}

% --- TeycoWeb ---
\item{\cventry{10/2014--10/2017}{Sviluppatore Web}{TeycoWeb S.r.l.}{Palermo, Italia}{}{\vspace{1pt}
\begin{itemize}
    \item Sviluppo e manutenzione evolutiva di back-office e servizi web con ASP.NET, C\#, SQL Server, XSLT e jQuery, a supporto dell’operatività di un provider di scommesse autorizzato.
    \item Realizzazione di componenti e infrastrutture applicative per Lotto Online, 10eLotto e Bet Exchange, integrando servizi e flussi informativi con partner esterni.
    \item Contributo alla continuità operativa di sistemi business-critical attraverso manutenzione, analisi dei problemi e ottimizzazione database.
\end{itemize}}}

\vspace{3pt}

% --- Leonardo ---
\item{\cventry{04/2007--09/2014}{Sviluppatore Web}{Leonardo Service Provider S.p.A.}{Palermo, Italia}{}{\vspace{1pt}
\begin{itemize}
    \item Progettazione di servizi interoperabili Java/.NET e database SQL Server; il servizio Gratta \& Vinci online ha superato il collaudo con throughput di 19 transazioni al secondo.
    \item Sviluppo di back-office e infrastrutture di servizio con ASP.NET, C\#, SQL Server e XML/XSLT; realizzazione di integrazioni per il protocollo ippico PSR TRIS - BIG MATCH.
\end{itemize}}}

\vspace{3pt}

% --- Older ---
\item{\cventry{2001--2007}{Sviluppatore software / System Engineer}{Esperienze precedenti selezionate}{Palermo, Italia}{}{\vspace{1pt}
\begin{itemize}
    \item Sviluppo di applicazioni ASP.NET, VB.NET/VB6 e SQL Server/MSDE per servizi pubblici, aziende e sistemi gestionali distribuiti.
    \item Realizzazione di un’applicazione client/server SOAP/XML per l’aggiornamento dati di circa 200 BetShop tra Serbia, Bosnia e Montenegro.
    \item Gestione di infrastrutture Windows Server, Active Directory, DHCP/DNS e misure di sicurezza.
\end{itemize}}}

\end{itemize}

% ============================================================
%     PROJECTS
% ============================================================

\section{Progetti Tecnici Selezionati}
\vspace{1pt}
\begin{itemize}

\item \textbf{CSharp.Essentials / LoggerHelper}: Progettazione di un'architettura modulare a plugin con registrazione ModuleInitializer, ottimizzazione del middleware HTTP con ArrayPool<char> e routing O(1) precomputato, e politiche di resilienza HTTP.

\item \textbf{FlowScheduler}: Client LLM resiliente con fallback/cooldown thread-safe e guardrail SQL read-only deterministici per bloccare DDL/DML generati da agenti AI.

\item \textbf{Server MCP .NET}: Esposizione di telemetria e diagnostica ad agenti AI tramite Model Context Protocol con Dependency Inversion.

\item \textbf{Pipeline real-time BetFair}: Ingestione dati streaming con System.Threading.Channels bounded e backpressure, e gestione scoped in worker a lungo ciclo tramite IServiceScopeFactory.

\item \textbf{Competenze AI e MCP}: Sviluppo di telemetria basata su Model Context Protocol e guardrail deterministici.

\end{itemize}

% ============================================================
%     EDUCATION
% ============================================================

\section{Formazione}
\vspace{1pt}
\begin{itemize}

\item{\cventry{1991--1998}{Laurea (Vecchio Ordinamento) in Matematica}{Università degli Studi di Palermo}{Palermo, Italia}{}{Conseguita il 18 marzo 1998 con votazione 100/110. Tesi: ``Metodi numerici e applicazioni matematiche.''}}

\end{itemize}

% ============================================================
%     LANGUAGES & CERTIFICATIONS
% ============================================================

\section{Lingue e Certificazioni}
\vspace{1pt}
\begin{itemize}
\item \textbf{Lingue}: Italiano (madrelingua), Inglese (livello B2), Spagnolo (livello A2).
\item \textbf{Certificazioni}: Certificazioni Microsoft (Windows 2000 Professional e Server, Network Infrastructure, Directory Services e Security).
\end{itemize}

% ============================================================
%     REFERENCES
% ============================================================

\section{Referenze}
\vspace{1pt}
\begin{itemize}
\item Disponibili su richiesta.
\end{itemize}

\end{document}
